%============================================================
\chapter{Estructura completa del proyecto}\label{ap:estructura}
%============================================================

\begin{lstlisting}[numbers=none,caption={Archivos del proyecto BiblioUTEQ},label={lst:arbol}]
biblio/
|-- pom.xml                      (proyecto padre: Spring Boot 4.1.2, Java 27, dos modulos)
|-- mvnw, mvnw.cmd, .mvn/        (Maven Wrapper: Maven 3.9.16)
|-- compose.yaml                 (tres niveles: web, api, db)
|-- README.md
|-- .github/workflows/ci.yml     (integracion continua)
|-- biblio-api/                  (nivel 2: servicios web, negocio, dominio y acceso a datos)
|   |-- Dockerfile
|   |-- pom.xml
|   `-- src/
|       |-- main/java/ec/edu/uteq/biblio/api/
|       |   |-- BiblioApiApplication.java
|       |   |-- config/          ConfiguracionReloj, ConfiguracionOpenApi
|       |   |-- web/             LibroControlador, SocioControlador, PrestamoControlador,
|       |   |                    ManejadorErrores
|       |   |-- web/dto/         LibroSolicitud, LibroRespuesta, SocioRespuesta,
|       |   |                    PrestamoSolicitud, PrestamoRespuesta
|       |   |-- negocio/         CatalogoServicio, PrestamoServicio, ReglasPrestamo,
|       |   |                    NotificadorPrestamos, RecursoNoEncontradoException
|       |   |-- dominio/         Libro, Socio, Prestamo, EstadoPrestamo,
|       |   |                    PrestamoRegistradoEvento, ReglaNegocioException
|       |   `-- datos/           LibroRepositorio, SocioRepositorio, PrestamoRepositorio
|       |-- main/resources/
|       |   |-- application.yaml
|       |   `-- db/migration/    V1__esquema_inicial.sql, V2__datos_iniciales.sql
|       `-- test/java/ec/edu/uteq/biblio/api/
|           |-- ContenedoresPrueba.java
|           |-- negocio/         PrestamoServicioTest
|           |-- web/             PrestamoControladorTest
|           |-- datos/           PrestamoRepositorioTest
|           `-- integracion/     BiblioApiIntegracionTest
`-- biblio-web/                  (nivel 1: presentacion)
    |-- Dockerfile
    |-- pom.xml
    `-- src/
        |-- main/java/ec/edu/uteq/biblio/web/
        |   |-- BiblioWebApplication.java
        |   |-- cliente/         BiblioApiCliente, Modelos
        |   |-- config/          ClienteApiConfig
        |   `-- controlador/     CatalogoWebControlador, PrestamoWebControlador,
        |                        ManejadorFallosApi
        |-- main/resources/
        |   |-- application.yaml
        |   |-- templates/       fragmentos, catalogo, prestamos, servicio-no-disponible
        |   `-- static/css/      estilos.css
        `-- test/java/ec/edu/uteq/biblio/web/
            `-- CatalogoWebControladorTest
\end{lstlisting}

%============================================================
\chapter{Comandos de referencia}\label{ap:comandos}
%============================================================

\begin{xltabular}{\linewidth}{@{}L{7.4cm}Y@{}}
	\caption{Comandos frecuentes en PowerShell}\label{tab:comandos}\\
	\toprule
	\rowcolor{uteqBlue}\hdr{Comando} & \hdr{Para qué sirve} \\
	\midrule
	\endfirsthead
	\toprule
	\rowcolor{uteqBlue}\hdr{Comando} & \hdr{Para qué sirve} \\
	\midrule
	\endhead
	\texttt{.\textbackslash mvnw.cmd verify} & Compilar y ejecutar todas las pruebas \\
	\texttt{.\textbackslash mvnw.cmd -DskipTests package} & Generar los \texttt{.jar} sin ejecutar pruebas \\
	\texttt{.\textbackslash mvnw.cmd -pl biblio-api test} & Ejecutar solo las pruebas del API \\
	\texttt{docker compose up -d db} & Levantar solo PostgreSQL para trabajar desde IntelliJ IDEA \\
	\texttt{docker compose up -d --build} & Construir las imágenes y levantar los tres niveles \\
	\texttt{docker compose ps} & Ver el estado de los contenedores \\
	\texttt{docker compose logs -f api} & Seguir los registros del API \\
	\texttt{docker compose stop api} & Detener un nivel para provocar una falla parcial \\
	\texttt{docker compose down} & Detener y eliminar los contenedores (conserva los datos) \\
	\texttt{docker compose down -v} & Eliminar también el volumen de datos \\
	\texttt{docker compose exec db psql -U biblio -d biblio} & Abrir una consola SQL en PostgreSQL \\
	\texttt{Invoke-RestMethod http://localhost:8080/api/v1/libros} & Consultar el API desde PowerShell \\
	\bottomrule
\end{xltabular}
